\chapter{Uchwyty Cassona i zanurzanie dysków}
\label{ch:uchwyty-cassona}
\index{uchwyt Cassona}
\index{zanurzanie dysków!w wymiarze cztery}

W rozdziale~\ref{ch:freedman-formy} obliczyliśmy formy
przecięcia i zobaczyliśmy, dlaczego gładki dysk Whitneya
może mieć nowe przecięcia w czterech wymiarach.
Teraz zastąpimy pojedynczą próbę usunięcia tych przecięć
ciągiem coraz wyższych pięter. Skończone piętra i ich
algebraiczną własność przeanalizujemy bezpośrednio.
Przejście od nieskończonej wieży do lokalnie płaskiego
zanurzenia jest głębokim twierdzeniem topologicznym;
podamy dokładnie jego hipotezy i źródła.

\section{Od zanurzonego dysku do pierwszego piętra}
\label{sec:casson-kinky}

Niech $M^4$ będzie zorientowaną gładką rozmaitością
z brzegiem, a
$f:(D^2,S^1)\to(M,\partial M)$ właściwym gładkim
zanurzeniem: $f(\operatorname{int}D^2)\subset
\operatorname{int}M$, zaś $f|_{S^1}$ jest zanurzeniem.
Zakładamy, że samoprzecięcia w środku są poprzecznymi
punktami podwójnymi i że przy brzegu $f$ ma zanurzony
kołnierz. Wybrane obramowanie tego kołnierza
ustala pas przyczepienia $S^1\times D^2$; jego
zgodność z kanonicznym obramowaniem następnego
modelu jest osobnym warunkiem, a nie skutkiem
samego istnienia zanurzenia.
Każdy punkt podwójny ma znak $+1$ albo $-1$,
zależnie od zgodności orientacji dwóch płatów z
orientacją $M$.

\begin{definicja}[Uchwyt z załamaniami]
\label{def:casson-kinky}
\emph{Uchwytem z załamaniami} odpowiadającym
$k$ punktom podwójnym nazywamy parę
$(K,\partial_-K)$ otrzymaną ze standardowego
uchwytu $D^2\times D^2$ przez $k$ samoplumbowań
poza obszarem przyczepienia
$\partial_-K=S^1\times D^2$. Rdzeń
$D^2\times\{0\}$ przechodzi w zanurzony dysk,
a znaki plumbowań są znakami jego punktów
podwójnych. Przy każdym takim punkcie wyróżniamy
w $\partial K\setminus\operatorname{int}\partial_-K$
okrąg przechodzący raz przez powstały $1$-uchwyt,
z \emph{zerowym obramowaniem} wyznaczonym przez
model plumbowania.
\end{definicja}

Lokalnie samoplumbowanie identyfikuje dwa małe,
rozłączne kawałki $D^2\times D^2$ leżące nad
różnymi poddyskami rdzenia, zamieniając rolę
czynnika rdzeniowego i normalnego. Ich rdzenie
spotykają się po identyfikacji w jednym punkcie.
Wybór orientacji identyfikacji ustala znak tego
punktu. Obramowanie wyróżnionego okręgu
pochodzi z tego modelu, więc nie trzeba go
wybierać ponownie przy budowie kolejnego piętra.
Równoważnie jest to takie obramowanie, przy którym
dołączenie zwykłego uchwytu $2$ do każdej z tych
pętli odzyskuje standardowy uchwyt $D^2\times D^2$
z początkowym obszarem przyczepienia.

Intuicja jest następująca: zamiast żądać, aby
zanurzony rdzeń od razu stał się zanurzony, zapamiętujemy
każde jego samoprzecięcie jako nową pętlę. Do
wyróżnionych pętli dołączymy potem następne uchwyty.
W przypadku jednego dodatniego punktu podwójnego
otrzymujemy jeden generator $\pi_1(K)\cong\Z$;
przy jednym ujemnym punkcie grupa abstrakcyjnej
przestrzeni jest taka sama, lecz obramowana para
$(K,\partial_-K)$ zachowuje znak plumbowania.
Model lokalny i obramowanie opisuje
\href{https://projecteuclid.org/journals/journal-of-differential-geometry/volume-17/issue-3/The-topology-of-four-dimensional-manifolds/10.4310/jdg/1214437136.full}
{Freedman, \emph{The topology of four-dimensional manifolds}
(1982), \S\,2}.

\begin{stwierdzenie}[Pętle pierwszego piętra]
\label{prop:casson-first-pi1}
Jeśli $K$ ma $k$ załamań, to $K$ ma typ homotopijny
sumy klinowej $k$ okręgów. Wyróżnione okręgi
odpowiadają wolnym generatorom $\pi_1(K)$.
\end{stwierdzenie}
\begin{proof}
W lokalnym rachunku uchwytów każde
samoplumbowanie dodaje jeden $1$-uchwyt; rdzeń starego
uchwytu $2$ nadal daje relację wynikającą z jego
okręgu przyczepienia. Okrąg ten jest brzegiem
zanurzonego rdzenia, więc relacja go dotycząca
nie wiąże nowych pętli przechodzących przez
$1$-uchwyty. Po zapomnieniu obszaru przyczepienia
rachunek uchwytów daje ciało z jednym $0$-uchwytem
i $k$ uchwytami $1$, a więc typ homotopijny sumy
klinowej $k$ okręgów.
Wyróżniony okrąg przy danym plumbowaniu przechodzi
raz przez odpowiadający mu $1$-uchwyt, a przez
pozostałe nie przechodzi, toteż reprezentuje
odpowiedni generator. Dla $k=0$ otrzymujemy
standardowy $D^2\times D^2$ i pustą sumę klinową.
Opis użytego rachunku znajduje się również w
cytowanym wyżej \S\,2 pracy Freedmana.
\end{proof}

\section{Wieże i drzewa znakowane}
\label{sec:casson-tower}

\begin{definicja}[Wieża i uchwyt Cassona]
\label{def:casson-handle}
Niech $T_1$ będzie uchwytem z załamaniami z co
najmniej jednym punktem podwójnym. Aby otrzymać
$T_{n+1}$, do każdej wyróżnionej pętli pochodzącej
z załamania na najwyższym piętrze $T_n$ dołączamy
nowy uchwyt z załamaniami, używając jej zerowego
obramowania. Każdy dołączany uchwyt ma skończoną,
niezerową liczbę załamań. Tak uzyskane $T_n$ nazywamy
\emph{wieżą $n$-piętrową}; odpowiednio otwarta suma
rosnąca $CH=\bigcup_{n\geq1}T_n$, z zachowaniem
początkowego obszaru przyczepienia, jest
\emph{uchwytem Cassona}. Przy sumie rozumiemy
standardowe przycięcie bocznych kołnierzy pięter,
tak aby $CH$ był gładką rozmaitością z wyróżnionym
otwartym obszarem przyczepienia
$\partial_-CH\cong S^1\times\operatorname{int}D^2$.
\end{definicja}

Każdy nowy rdzeń dysku wypełnia jedną pętlę
poprzedniego piętra, lecz może mieć własne punkty
podwójne. Dlatego konstrukcja trwa bez końca.
Warunek zerowego obramowania nie jest ozdobą:
utrzymuje właściwy typ końca uchwytu. Jego rolę
wyjaśnia Freedman, \S\,2, bezpośrednio po definicji
uchwytu Cassona.

Wieżę koduje zakorzenione drzewo lokalnie skończone
w sensie Definicji~\ref{def:drzewo-grafowe}:
wierzchołek to uchwyt z załamaniami, a każda krawędź
wychodząca z niego odpowiada jednemu punktowi
podwójnemu jego rdzenia. Krawędź nosi znak tego
punktu. Nie dopuszczamy liści, bo każdy punkt
podwójny otrzymuje następne piętro. Taki rysunek
koduje konstrukcję, ale sam nie jest twierdzeniem
o klasyfikacji gładkiej uchwytów Cassona.

\begin{figure}[htbp]
\centering
\begin{tikzpicture}[>=Latex,font=\small,
 cnode/.style={circle,draw=manifoldblue!80!black,
 fill=manifoldblue!12,minimum size=9mm,inner sep=1pt}]
 \node[cnode] (r) at (0,2.75) {$K_0$};
 \node[cnode] (a) at (-2,1.35) {$K_+$};
 \node[cnode] (b) at (2,1.35) {$K_-$};
 \node[cnode] (c) at (-2,0) {$K_{++}$};
 \node[cnode] (d) at (1,0) {$K_{-+}$};
 \node[cnode] (e) at (3,0) {$K_{--}$};
 \draw[manifoldblue!80!black,thick]
   (r)--node[above left] {$+$} (a)
   (r)--node[above right] {$-$} (b)
   (a)--node[right] {$+$} (c)
   (b)--node[above left] {$+$} (d)
   (b)--node[above right] {$-$} (e);
 \foreach \x in {c,d,e}{
   \draw[manifoldblue!65!black,densely dotted,thick]
      (\x.south)--++(0,-.45);
 }
 \node[manifoldblue!75!black] at (-2,-.8) {$\vdots$};
 \node[manifoldblue!75!black] at (1,-.8) {$\vdots$};
 \node[manifoldblue!75!black] at (3,-.8) {$\vdots$};
\end{tikzpicture}
\caption{Pierwsze trzy poziomy drzewa. $K_0$ ma
dwa załamania o znakach $+$ i $-$; dlatego w drugim
piętrze są dwa uchwyty. Ich załamania wyznaczają
trzecie piętro. Kropki oznaczają dalsze poziomy,
a nie końcowe wierzchołki.}
\label{fig:casson-signed-tree}
\end{figure}

\begin{przyklad}[Jedna gałąź]
\label{prz:casson-ray}
Jeśli każdy uchwyt ma dokładnie jedno dodatnie
załamanie, drzewo jest nieskończonym promieniem
ze wszystkimi znakami $+$. Każde piętro $T_n$
ma na szczycie jedną pętlę. Zmiana jednego znaku
nie zmienia tego prostego rachunku grupy podstawowej,
choć zmienia dane gładkich plumbowań.
\end{przyklad}

\begin{stwierdzenie}[Co dają skończone piętra]
\label{prop:casson-tower-pi1}
Niech $r_n$ oznacza liczbę punktów podwójnych
na $n$-tym piętrze. Wtedy
$\pi_1(T_n)$ jest grupą wolną na $r_n$ generatorach,
a homomorfizm
$\pi_1(T_n)\to\pi_1(T_{n+1})$ jest zerowy.
W szczególności $\pi_1(CH)=0$.
\end{stwierdzenie}
\begin{proof}
Dla $T_1$ jest to
Stwierdzenie~\ref{prop:casson-first-pi1}.
Załóżmy, że generatory $\pi_1(T_n)$ są pętlami
odpowiadającymi załamaniom najwyższego piętra.
Okrąg przyczepienia każdego nowego uchwytu
reprezentuje dokładnie jeden z tych generatorów.
W nowym uchwycie ten okrąg jest brzegiem jego
zanurzonego rdzenia, a więc staje się
nullhomotopijny. Twierdzenie van Kampena usuwa
wszystkie stare generatory. Pozostają wolne
generatory utworzone przez załamania nowych
uchwytów; ich liczba wynosi $r_{n+1}$.
To kończy indukcję. Każda pętla w $CH$ ma zwarty
obraz, mieści się więc w pewnym $T_n$ i staje się
nullhomotopijna w $T_{n+1}$.
\end{proof}

Ten rachunek nie stwierdza, że $CH$ jest standardowym
otwartym uchwytem $2$. Sama grupa podstawowa nie
rozpoznaje zachowania przy końcu niezwartym ani
lokalnej płaskości ewentualnego dysku granicznego.

\section{Dyski Whitneya i dane w pierścieniu grupowym}
\label{sec:casson-whitney}

Na Rysunku~\ref{fig:freedman-whitney} widzieliśmy
parę punktów przecięcia o znakach $+$ i $-$.
Jeśli odpowiedni okrąg Whitneya jest
nullhomotopijny, można go wypełnić zanurzonym
dyskiem. Ruch Whitneya usuwa wyjściową parę,
gdy wnętrze tego dysku jest czyste i obramowanie
jest właściwe. W wymiarze cztery wnętrze może
przeciąć powierzchnie, a sam dysk może przeciąć
siebie. Przeniesienie problemu na następny dysk
jest ideą wspólną wieżom Cassona i konstrukcjom
z dyskami Whitneya. Są to różne modele geometryczne:
w klasycznej wieży Cassona następne rdzenie
dołącza się do pętli pojedynczych załamań, nie
po prostu do każdej pary punktów Whitneya.

Przy nietrywialnej $\pi_1(M)$ same znaki przecięć
nie wystarczają nawet do znalezienia okręgu
Whitneya, który można wypełnić. Dla dwóch
zorientowanych, obramowanych zanurzonych dysków
lub sfer $A,B$ wybieramy ścieżki od ustalonego
punktu bazowego do ich obrazów. Punkt
$p\in A\cap B$ daje znak $\varepsilon_p\in\{+1,-1\}$
i element $g_p\in\pi_1(M)$: idziemy do $A$,
po $A$ do $p$, a następnie wracamy po $B$.
Otrzymujemy liczbę przecięcia z wartościami w
pierścieniu grupowym
\begin{equation}\label{eq:casson-lambda}
 \lambda(A,B)=\sum_{p\in A\cap B}
             \varepsilon_p g_p\in\Z[\pi_1(M)].
\end{equation}
Przy samoprzecięciu wybór kolejności dwóch płatów
zmienia $g$ na $g^{-1}$. Dlatego obramowana
liczba samoprzecięcia $\mu(B)$ należy do
addytywnego ilorazu
$\Z[\pi_1(M)]/\langle g-g^{-1}:g\in\pi_1(M)\rangle$.
Zależność od orientacji i obramowania oraz definicję
ogólną podają
\href{https://archive.mpim-bonn.mpg.de/4789/2/FreedmanQuinn-TopologyOf4Manifolds-Reformatted2013.pdf}
{Freedman--Quinn, \emph{Topology of 4-Manifolds},
\S\,1.7--1.8}. W przypadku $\pi_1(M)=0$
wartości~\eqref{eq:casson-lambda} są po prostu
całkowitymi liczbami przecięcia.

\begin{stwierdzenie}[Algebraiczne pary punktów]
\label{prop:casson-pairing}
Jeżeli $\lambda(A,B)=0$, punkty $A\cap B$ można
sparować tak, aby w każdej parze znaki były
przeciwne, elementy $g_p$ jednakowe, a okrąg
Whitneya był nullhomotopijny w $M$. Jeżeli
$\lambda(A,B)=1$, po takim parowaniu pozostaje
jeden dodatni punkt o etykiecie $1\in\pi_1(M)$.
\end{stwierdzenie}
\begin{proof}
Pierścień grupowy jest wolną grupą abelową z bazą
$\pi_1(M)$. Zatem współczynnik każdego $g$
w sumie~\eqref{eq:casson-lambda} znika osobno,
z wyjątkiem współczynnika $1$ w drugim przypadku.
Kojarzymy dodatnie i ujemne punkty o tej samej
etykiecie. Dla takiej pary wybieramy łuk od
pierwszego punktu do drugiego na $A$ i łuk
powrotny na $B$. Definicja $g_p$ pokazuje,
że klasa otrzymanej pętli jest iloczynem
$g_p g_q^{-1}=1$. Pętla jest więc nullhomotopijna.
Jej wypełnienie można doprowadzić do położenia
ogólnego jako \emph{zanurzony} dysk Whitneya;
ten ostatni krok nie usuwa jego przecięć wewnętrznych.
\end{proof}

Analogiczna implikacja dla $\mu(B)=0$ wymaga
uwzględnienia zamiany płatów i opisanej relacji
$g\sim g^{-1}$; zob.
Freedman--Quinn, \S\,1.7. Obramowania też
nie można pominąć: skręcenie brzegu dysku
Whitneya może poprawić jego obramowanie,
ale zwykle wytwarza nowe przecięcie z jedną
z powierzchni. Ta algebraiczna redukcja daje
więc \emph{zanurzone} dyski; ich topologiczne
zanurzenie jest następnym, odrębnym krokiem.

\section{Twierdzenie o zanurzaniu dysku}
\label{sec:casson-disk-theorem}

Grupę $G$ nazywamy \emph{wielostopniowo skończoną
lub cykliczną}, jeśli ma skończony ciąg
$1=G_0\triangleleft G_1\triangleleft\cdots
\triangleleft G_m=G$, w którym każdy iloraz
$G_{i+1}/G_i$ jest skończony albo cykliczny.
W szczególności należą tu grupa trywialna,
grupy skończone i $\Z^k$. Jest to konkretna
klasa grup użyta w poniższej wersji twierdzenia.

\begin{twierdzenie}[Freedman--Quinn: zanurzanie dysków;
wynik zewnętrzny]
\label{thm:casson-disk-embedding}
Niech $M^4$ będzie zwartą, spójną, zorientowaną gładką
rozmaitością z brzegiem, a
$A_i:(D^2,S^1)\to(M,\partial M)$,
$1\leq i\leq r$, skończoną rodziną właściwie
zanurzonych obramowanych dysków. Ich brzegi
są parami rozłącznie zanurzone wraz z kołnierzami.
Załóżmy, że istnieją obramowane zanurzone sfery
$B_j:S^2\to\operatorname{int}M$ takie, że w
$\Z[\pi_1(M)]$ zachodzą
\begin{equation}\label{eq:casson-duals}
 \lambda(A_i,B_j)=\delta_{ij},\qquad
 \lambda(B_i,B_j)=0,
\end{equation}
a w grupie samoprzecięć zachodzi $\mu(B_j)=0$
dla każdego $j$. Jeśli $\pi_1(M)$ jest
wielostopniowo skończona lub cykliczna, to w $M$
istnieją parami rozłączne, lokalnie płasko
\emph{topologicznie zanurzone} dyski o tych samych
obramowanych brzegach. Mają one zanurzone sfery
poprzeczne, z których każda przecina odpowiedni
dysk w jednym punkcie i jest rozłączna z pozostałymi
dyskami.
\end{twierdzenie}

Sfery $B_j$ są \emph{algebraicznie dualne} do
dysków: pierwsza równość w~\eqref{eq:casson-duals}
oznacza jeden punkt przecięcia po skasowaniu
algebraicznych par. Pozostałe warunki kontrolują
przecięcia samych sfer. Obramowanie oznacza tu
trywializację dwuwymiarowej wiązki normalnej,
zgodną z zadanym obramowaniem brzegów dysków.
Twierdzenie podajemy za
\href{https://archive.mpim-bonn.mpg.de/4789/2/FreedmanQuinn-TopologyOf4Manifolds-Reformatted2013.pdf}
{Freedman--Quinn, \emph{Topology of 4-Manifolds},
Twierdzenie 5.2, s.\ 85--86}. Wersja dla gładkiego
$M$ unika dodatkowego przejścia od rozmaitości
topologicznej do gładkich lokalnych modeli
zanurzeń; to przejście autorzy omawiają
w rozdziale 9. Konkluzja pozostaje topologiczna.
Twierdzenie samo nie mówi, że wyjściowe $A_i$
są homotopijne regularnie do uzyskanych zanurzeń.

W konkluzji słowo \emph{poprzeczne} ma znaczenie
geometryczne: każda uzyskana sfera przecina swój
dysk w dokładnie jednym punkcie, a pozostałe dyski
omija. Samo równanie
$\lambda(A_i,B_j)=\delta_{ij}$ jest słabsze, bo
dopuszcza dodatkowe pary przecięć. W oryginalnym
dowodzie Freedmana--Quinna przejście od dualności
algebraicznej do takich sfer zawierało lukę. Domyka
ją \href{https://doi.org/10.1007/s00029-025-01069-y}
{Powell--Ray--Teichner, \emph{The 4-dimensional disc
embedding theorem and dual spheres} (2025),
Twierdzenie A i \S\,5}. Twierdzenie A obejmuje
przytoczoną wyżej wersję dla gładkiego $M$ i grup
wielostopniowo skończonych lub cyklicznych; zachowuje
też klasy homotopii sfer dualnych. Samo istnienie
zanurzonych dysków wynikało już z wcześniejszej
części argumentu Freedmana--Quinna. Nowe źródło
uzasadnia dodatkową, używaną dalej konkluzję
o \emph{geometrycznie} dualnych sferach.

\begin{stwierdzenie}[Co zapewnia geometryczna dualność]
\label{prop:casson-complement-pi1}
Niech $D_1,\ldots,D_r$ będą parami rozłącznymi,
lokalnie płasko zanurzonymi właściwymi dyskami w
zwartej spójnej rozmaitości $M^4$ z brzegiem.
Jeżeli mają geometrycznie
dualne zanurzone sfery, to inkluzja dopełnienia
indukuje izomorfizm
\[
 \pi_1\bigl(M\setminus\textstyle\bigcup_i D_i\bigr)
       \xrightarrow{\ \cong\ }\pi_1(M).
\]
\end{stwierdzenie}
\begin{proof}
Przez położenie ogólne każdą pętlę w $M$ można
odsunąć od dysków, więc homomorfizm jest
surjektywny. Niech pętla w dopełnieniu ogranicza
dysk odwzorowany w $M$. Korzystamy tu z
topologicznej transwersalności dla lokalnie
płaskich podrozmaitości oraz z reprezentowania
klas map przez generyczne immersje,
\href{https://arxiv.org/abs/2006.05209}
{Powell--Ray--Teichner, Twierdzenia 6.4 i 1.5}.
Po uczynieniu odwzorowania poprzecznym do $D_i$
przeciwobrazy
dysków są skończonym zbiorem punktów. Usuwamy
małe kółka wokół tych punktów. Pozostała część
dysku leży w dopełnieniu, a każdy nowy okrąg
brzegowy jest meridianem jednego z $D_i$.
Sfera dualna do $D_i$, po usunięciu małego krążka
otaczającego jej jedyny punkt przecięcia z $D_i$,
daje zanurzony dysk w dopełnieniu, którego brzeg
jest takim meridianem. Wszystkie meridiany są
zatem nullhomotopijne w dopełnieniu, także po
sprzężeniu ścieżkami potrzebnymi do zmiany punktu
bazowego. Doklejenie tych zanurzonych dysków
wypełnia wyjściową pętlę w dopełnieniu. Homomorfizm
jest więc także injektywny.
\end{proof}

Ten prosty skutek wyjaśnia, dlaczego nie wystarcza
zastąpić w konkluzji twierdzenia dualności
geometrycznej dualnością algebraiczną: w drugim
przypadku sfera może mieć kilka punktów przecięcia
z dyskiem, a powyższe wypełnienie meridianu nie
powstaje. Jest to również powód, dla którego
\href{https://doi.org/10.1007/s00029-025-01069-y}
{Powell--Ray--Teichner, Uwaga 1.2} podkreślają
kontrolę grupy podstawowej dopełnienia.

\begin{przyklad}[Standardowa para dualna]
\label{prz:casson-dual-product}
Weźmy $M=S^2\times D^2$,
$A=\{p\}\times D^2$ i $B=S^2\times\{0\}$.
Wtedy $\pi_1(M)=0$, dysk i sfera przecinają się
w jednym dodatnim punkcie, $\lambda(A,B)=1$,
a $\mu(B)=0$. Normalne kierunki do obu czynników
dają wymagane obramowania. Twierdzenie obejmuje
tę parę, chociaż tutaj $A$ jest już zanurzony.
Jeżeli zanurzenie $A$ zmienimy we wnętrzu
lokalnym ruchem palcowym, jego brzeg i dualność
algebraiczna nadal pozostają takie same;
nowe punkty podwójne pokazują, jak wyglądają
dane przed zastosowaniem twierdzenia.
\end{przyklad}

\begin{wniosek}[Wersja z homotopią regularną;
wynik zewnętrzny]
\label{cor:casson-regular-homotopy}
Przy założeniach Twierdzenia~\ref{thm:casson-disk-embedding}
załóżmy dodatkowo, że wszystkie algebraiczne
przecięcia wzajemne $A_i$ oraz ich liczby
samoprzecięcia $\mu(A_i)$ są zerowe.
Wtedy rodzinę $A_i$ można przez homotopię
regularną względnie brzegów doprowadzić do
parami rozłącznych topologicznie zanurzonych
dysków z dualnymi sferami.
\end{wniosek}

To jest \href{https://archive.mpim-bonn.mpg.de/4789/2/FreedmanQuinn-TopologyOf4Manifolds-Reformatted2013.pdf}
{Freedman--Quinn, Wniosek 5.2.1, s.\ 87--88}.
Dodatkowych równości dotyczących $A_i$ nie
ukrywamy w Twierdzeniu~\ref{thm:casson-disk-embedding}:
tam potrzebne jest jedynie zachowanie
obramowanych brzegów.

Warto rozdzielić pośrednie kroki tego wniosku.
Równości $\lambda(A_i,A_j)=0$ i $\mu(A_i)=0$
parują algebraicznie przecięcia dysków i pozwalają
wybrać \emph{zanurzone}, odpowiednio obramowane
dyski Whitneya. Najpierw zanurzone ruchy Whitneya
redukują dodatkowe przecięcia $A_i$ ze sferami
algebraicznie dualnymi. Następnie z torusów
opasujących punkty podwójne oraz sfer dualnych
buduje się sfery algebraicznie dualne do samych
dysków Whitneya. Do \emph{nich} stosuje się
Twierdzenie~\ref{thm:casson-disk-embedding},
otrzymując rozłączne zanurzone dyski Whitneya.
Prowadzone po nich ruchy Whitneya usuwają
przecięcia $A_i$; ponieważ są to ruchy wewnątrz,
zachowują obramowane brzegi i dają wskazaną
homotopię regularną. Konstrukcja sfer dualnych
do dysków Whitneya jest częścią zewnętrznego
dowodu Freedmana--Quinna, s.\ 87--88; nie wynika
z samego parowania w pierścieniu grupowym.

Zakres grup można opisać szerzej. Grupa skończenie
generowana ma \emph{wzrost podwykładniczy}, jeśli
dla dowolnego skończonego zbioru generatorów $S$
liczba elementów o długości co najwyżej $n$
spełnia $\lim_{n\to\infty}|B_S(n)|^{1/n}=1$.
Ta własność obejmuje w szczególności grupy
o wzroście wielomianowym. Dla grup podstawowych
o wzroście podwykładniczym tę samą konkluzję
zanurzania otrzymuje się dzięki ulepszeniu
kontroli pętli punktów podwójnych; zob.
\href{https://math.berkeley.edu/~teichner/Papers/Freedman1.pdf}
{Freedman--Teichner, \emph{4-Manifold topology I:
Subexponential groups} (1995)} oraz
\href{https://uva.theopenscholar.com/files/slava-krushkal/files/subexponential_groups.pdf}
{Krushkal--Quinn, \emph{Subexponential groups in
4-manifold topology} (2000)}. Ostatnia praca
podaje poprawiony argument i wskazuje błąd
w pierwotnym dowodzie z 1995 roku. Nie
rozszerzamy stąd twierdzenia na dowolną grupę
podstawową. Jest również wersja z dowolną
$\pi_1(M)$ przy silnym \emph{geometrycznym}
założeniu $\pi_1$-zerowości poprzecznych
powierzchni z dyskami zamykającymi
(Freedman--Quinn, Twierdzenie 5.4).

\section{Gdzie kończy się redukcja geometryczna}
\label{sec:casson-boundary}

Możemy teraz uporządkować rzeczywiście wykonane
kroki. Zerowanie współczynników pierścienia
grupowego pozwala sparować punkty przecięcia
i znaleźć zanurzone dyski Whitneya
(Stwierdzenie~\ref{prop:casson-pairing}).
Samoplumbowanie daje uchwyt pierwszego
piętra z podpisanym znakiem każdego punktu
podwójnego. Dołączanie kolejnych pięter zabija
generatory grupy podstawowej powstałe na
poprzednim piętrze
(Stwierdzenie~\ref{prop:casson-tower-pi1}).
Żaden z tych skończonych rachunków nie produkuje
jeszcze dysku o czystym wnętrzu.

W dowodzie Twierdzenia~\ref{thm:casson-disk-embedding}
potrzeba ponadto skonstruować powierzchnie z
dyskami zamykającymi i dualne sfery, poprawić
grupę podstawową pętli punktów podwójnych,
otrzymać zbieżną wieżę oraz zanurzyć dysk
lokalnie płasko w jej otoczeniu. Są to zasadnicze
zewnętrzne części teorii Freedmana--Quinna:
\href{https://archive.mpim-bonn.mpg.de/4789/2/FreedmanQuinn-TopologyOf4Manifolds-Reformatted2013.pdf}
{\emph{Topology of 4-Manifolds}, rozdziały 2--5}.
Ostatnie przejście opiera się na kontroli
topologicznej przy nieskończonym końcu wieży;
nie wynika z samego van Kampena ani z formalnej
granicy ciągu ruchów Whitneya.

\begin{twierdzenie}[Freedman: standardowość topologiczna
uchwytu Cassona; wynik zewnętrzny]
\label{thm:casson-homeomorphism}
Każdy uchwyt Cassona wraz ze swym obszarem
przyczepienia jest homeomorficzny jako para ze
standardowym otwartym uchwytem $2$:
\[
 (CH,\partial_-CH)\cong
 (D^2\times\R^2,S^1\times\R^2).
\]
Homeomorfizm można dobrać zgodnie z ustalonym
obramowaniem obszaru przyczepienia.
\end{twierdzenie}

Jest to \href{https://projecteuclid.org/journals/journal-of-differential-geometry/volume-17/issue-3/The-topology-of-four-dimensional-manifolds/10.4310/jdg/1214437136.full}
{Freedman, \emph{The topology of four-dimensional
manifolds} (1982), Twierdzenie 1.1}.
Przyjmujemy je tutaj jako wynik zewnętrzny.
Jego dowód wymaga ponownego zanurzania wież,
kontroli zbiorów granicznych i argumentów
zmniejszania rozkładów topologicznych. Dopiero
ten wynik uzasadnia przejście od uchwytu Cassona
do standardowego otwartego uchwytu w kategorii
TOP. Nie daje analogicznego dyfeomorfizmu par
w kategorii gładkiej. W szczególności topologiczny
rdzeń $D^2\times\{0\}$ nie musi być gładkim
zanurzonym dyskiem w pierwotnym gładkim $M$.

Twierdzenie o zanurzaniu dysku jest geometrycznym
narzędziem stojącym za topologicznymi wnioskami
Freedmana przytoczonymi w
rozdziale~\ref{ch:freedman-formy}. Wnioski o
gładkich strukturach wymagają osobnych przeszkód;
same znakowane drzewa, wyzerowanie $\pi_1(CH)$
i lokalna płaskość topologiczna ich nie rozstrzygają.
